\chapter{Invitation to real analysis}

This part of the book is a bit out of touch with the other chapters. I named this whole series \enquote{Compact Calculus - Calculus for the practitioners,} because I envision it to be a gentle and compact introduction to calculus that is used by most practitioners. One could argue that real analysis isn't for the practitioners, as it is considered to be one of the hardest pure mathematics course that an undergraduate can take. It's dedicated to rigorously analyzing the behaviors of the reals, redefining many of the hand-wavy definitions that we've used, and proving even the most trivial properties. Some proofs are so abstract that most have trouble wrapping their heads around them. So, why did I decide to write about real analysis?

After reading many analysis book\footnote{Even though I don't even have the point of doing so; I do physics, not mathematics.}, I've developed a deeper appreciation for the nature of mathematics. It's remarkable that from just a few axioms, mathematics has managed to prosper into the vast interconnecting webs of theories that we know today. So, I'd like to take my appreciation and write them down on a book, where those who are interested can also appreciate the foundations of the reals like I did.

